% =============================================================================
% Anhang: Referenzlauf 930444 – Testdaten und exemplarische Traceability-Evidence
% =============================================================================

\section{Referenzlauf 930444: Testdaten und exemplarische Traceability-Evidence}
\label{app:traceability_930444}

Dieser Anhang dokumentiert die im Haupttext in
Abschnitt~\ref{sec:results_traceability} exemplarisch dargestellte
End-to-End-Traceability des Referenzprojekts~930444. Die Darstellung verfolgt
zwei Ziele: Zum einen werden die im Referenzlauf registrierten und für die
Evaluation aufbereiteten Eingangsdaten vollständig offengelegt. Zum anderen
werden für den exemplarisch gewählten Informationspfad
\emph{Microscope Workstation} die für die Rekonstruktion relevanten
persistierten Artefakte in Originalauszügen zusammengeführt.

Der vollständige Lauf erzeugte einen wesentlich größeren Evidence-, Decision-,
Provenance- und Traceability-Artefaktsatz. Eine vollständige Wiedergabe aller
maschinenlesbaren JSON-Artefakte würde den Anhang unverhältnismäßig erweitern.
Daher werden die kurzen Eingangsdokumente und der kompakte Validation Report
vollständig wiedergegeben. Umfangreiche Reports und Maschinenartefakte werden auf die für
den dargestellten Pfad relevanten Felder reduziert. Artefaktidentitäten,
Vorgängerbezüge, Reviewer, Rationale und Authority-Referenzen bleiben dabei
sichtbar.

Die Reviewer-Identität \texttt{MZ} bezeichnet das im Prototyp verwendete
Kürzel des Autors. Es leitet sich aus \emph{MoritZ} ab und wurde während des
Referenzlaufs für die expliziten Human-Review-, Model-Placement-,
Model-Quality-, Target-Formulation- und Release-Entscheidungen verwendet.

\subsection{Aufbereitete Engineering-Eingangsdaten des Referenzlaufs}
\label{app:traceability_930444_inputs}

Die im Referenzlauf verwendeten Dokumente basieren auf realen
Engineering-Inhalten von PreciPoint zum betrachteten
Remote-Collaboration-Kontext. Für den abschließenden Referenzlauf wurden die
Inhalte gegenüber den umfangreicheren formativen Testdaten weiter reduziert
und strukturiert. Ziel war es, den vollständigen Informations- und
Authority-Pfad anhand klar abgrenzbarer Engineering-Aussagen reproduzierbar
zu untersuchen und gleichzeitig Laufzeit und Tokenbedarf der wiederholten
LLM-Verarbeitung zu begrenzen. Die Dokumente sind damit keine unveränderten
Produktivdokumente, ihre fachliche Grundlage stammt jedoch aus realen
Engineering-Daten.

Für den Referenzlauf wurden sechs Quellen separat registriert. Vier Quellen
wurden als \texttt{engineering\_source} verarbeitet. Zwei kurze
Produktkontexte wurden als \texttt{context\_only} registriert und dienten dem
projektbezogenen Interpretations- und Project-Fit-Kontext.

\begin{table}[htbp]
    \centering
    \small
    \caption{Registrierte Quellen des Referenzprojekts 930444}
    \label{tab:app_930444_sources}
    \begin{tabularx}{\textwidth}{p{2.1cm}p{4.0cm}p{3.2cm}X}
        \hline
        \textbf{Source} & \textbf{Datei} & \textbf{Rolle} & \textbf{Dateipfad} \\
        \hline
        \texttt{SRC-000001} & \texttt{01\_product\_context.md} & \texttt{context\_only} & \path{safe_demo/01_product_context.md} \\
        \texttt{SRC-000002} & \texttt{02\_stakeholders.md} & \texttt{engineering\_source} & \path{safe_demo/02_stakeholders.md} \\
        \texttt{SRC-000003} & \texttt{03\_user\_workflow.md} & \texttt{engineering\_source} & \path{safe_demo/03_user_workflow.md} \\
        \texttt{SRC-000004} & \texttt{04\_system\_requirements.md} & \texttt{engineering\_source} & \path{safe_demo/04_system_requirements.md} \\
        \texttt{SRC-000005} & \texttt{05\_architecture\_notes.md} & \texttt{engineering\_source} & \path{safe_demo/05_architecture_notes.md} \\
        \texttt{SRC-000006} & \texttt{01\_product\_context.md} & \texttt{context\_only} & \path{safe_demo/live_minimal/01_product_context.md} \\
        \hline
    \end{tabularx}
\end{table}

Die zugehörigen SHA-256-Werte wurden mit den registrierten Source-Manifests
persistiert:

\begin{evidencebox}
SRC-000001  81aeec42e1a5971e8951312edb1a6fe4466e267b1f102e4e147f5ff262e51a86
SRC-000002  1138b4f0d55b161492fc0a4d33505965465d554c00590653808f89cac3bd6624
SRC-000003  cef50598657ae53b9f03db36567f50752fc03caf7685980bccdd79029c8df12e
SRC-000004  d7c7c73a8ebe0e2483b0b54c63758cd918159fa2af835e2b1a8e9b1a26c4ec26
SRC-000005  a4cd264fdbdd518999d4eb162d8d8a3e97a1df57f5f75ca53bb1ebfcbc4d0e77
SRC-000006  1590465ec7c01cffde3a9567b96cfafca4f9213d6f1a8198f031219365933faa
\end{evidencebox}

\paragraph{Produktkontext, \texttt{SRC-000001}.}
Originaldatei: \path{legacy/demo/safe_demo/01_product_context.md}.

\begin{evidencebox}
# Product Context

iO Stream supports remote collaboration in digital microscopy.

A Microscope Operator works locally and a Remote Expert participates remotely.
\end{evidencebox}

\paragraph{Stakeholder-Beschreibung, \texttt{SRC-000002}.}
Originaldatei: \path{legacy/demo/safe_demo/02_stakeholders.md}.

\begin{evidencebox}
# Stakeholders

The Microscope Operator is the local user.

The Remote Expert is the remote user.
\end{evidencebox}

\paragraph{User Workflow, \texttt{SRC-000003}.}
Originaldatei: \path{legacy/demo/safe_demo/03_user_workflow.md}.

\begin{evidencebox}
# User Workflow

1. The Microscope Operator starts the collaboration session.
2. The Remote Expert requests microscope control.
3. The Microscope Operator may take back control.
\end{evidencebox}

\paragraph{System Requirements, \texttt{SRC-000004}.}
Originaldatei: \path{legacy/demo/safe_demo/04_system_requirements.md}.

\begin{evidencebox}
# System Requirements

The system shall indicate the current controller.

The system shall disable remote control after connection loss.

The system shall return control to the Microscope Operator after connection loss.
\end{evidencebox}

\paragraph{Architecture Notes, \texttt{SRC-000005}.}
Originaldatei: \path{legacy/demo/safe_demo/05_architecture_notes.md}.
Diese Quelle enthält den im folgenden Traceability-Beispiel verfolgten
Gegenstand \emph{Microscope Workstation}.

\begin{evidencebox}
# Architecture Notes

The Microscope Workstation provides local session control.

The Remote Client provides remote viewing and control requests.
\end{evidencebox}

\paragraph{Minimaler Produktkontext, \texttt{SRC-000006}.}
Originaldatei:
\path{legacy/demo/safe_demo/live_minimal/01_product_context.md}.

\begin{evidencebox}
# Product Context

iO Stream supports remote collaboration in digital microscopy.
The main users are a Microscope Operator and a Remote Expert.
\end{evidencebox}

\subsection{Exemplarischer Traceability-Pfad: Microscope Workstation}
\label{app:traceability_930444_workstation}

Der folgende Nachweis verfolgt einen einzelnen Informationspfad vom
registrierten Source-Inhalt bis zum publizierten SysML-v2-Artefakt. Der Pfad
wurde ausgewählt, weil innerhalb desselben Gegenstands mehrere unterschiedliche
Verarbeitungszustände sichtbar sind: ein fehlgeschlagener Processing Attempt,
ein expliziter Retry, divergierende Persona-Klassifikationen, eine Human
Engineering Decision, Model Placement, eine nachgelagerte Model-Quality-
Entscheidung sowie die finale Output-Traceability.

\subsubsection{Source und fehlgeschlagener Processing Attempt}

\texttt{SRC-000005} wurde in \texttt{RUN-000001} als
\texttt{engineering\_source} verarbeitet. Der erste Attempt
\texttt{ATT-000001} wechselte aus dem Zustand \texttt{running} in
\texttt{failed}. Der Fehlerzustand wurde als eigenes Processing Event
persistiert.

\begin{evidencebox}
event_id:         EVT-000003
previous_state:   running
next_state:       failed
processing_stage: agentic_ingestion
event_type:       run_failed
attempt_id:       ATT-000001
reason_code:      shared_evidence_interpretation_failed
\end{evidencebox}

\noindent\textit{Originalartefakt:}
\texttt{runs/RUN-000001/events/EVT-000003.json}.

Für den fehlgeschlagenen Attempt blieb zugleich die rohe Antwort des
Literal-Interpreters erhalten. Der folgende Ausschnitt zeigt das im ersten
Interpretationseintrag zurückgegebene Feld \texttt{extracton\_rationale}.

\begin{evidencebox}
{
  "source_evidence_id": "EVD-000001",
  "interpreted_statement": "This source is labeled as architecture notes.",
  "information_type": "definition",
  "statement_modality": "descriptive",
  "epistemic_class": "explicit",
  "missing_evidence": null,
  "extracton_rationale": "The evidence explicitly names the source material as architecture notes, but provides no further engineering content.",
  "uncertainties": [
    "The note title is broad and does not specify a concrete system behavior or element."
  ]
}
\end{evidencebox}

\noindent\textit{Originalartefakt:}
\texttt{ATT-000001/.../agent\_legacy\_literal\_interpreter\_run\_01.json}.
Der zugehörige Parsing-Contract erwartet an dieser Stelle das Feld
\texttt{extraction\_rationale}. Der fehlerhafte Attempt wurde nicht als
nachgelagerter Engineering-Zustand weitergeführt. Stattdessen wurde mit
\texttt{ATT-000002} ein neuer Attempt begonnen. Die Event-Historie persistiert
diesen Übergang separat:

\begin{evidencebox}
event_id:       EVT-000004
previous_state: failed
next_state:     running
event_type:     retry_started
attempt_id:     ATT-000002
reason_code:    agentic_ingestion_retry_started
\end{evidencebox}

\subsubsection{Source-grounded Evidence und Canonical Subject}

Im erfolgreichen zweiten Attempt wurde die Source-Aussage als
\texttt{EVD-000002} mit Source- und Segmentbezug persistiert.

\begin{evidencebox}
source_evidence_id: EVD-000002
source_id:          SRC-000005
source_projection:  SP-000001
source_anchor:      SEG-000002 [0, 58]
source_excerpt:     The Microscope Workstation provides local session control.
\end{evidencebox}

\noindent\textit{Originalartefakt:}
\texttt{semantics/source\_evidence/EVD-000002.json}.

Die anschließende Subject Discovery führte den Gegenstand als gemeinsames
Canonical Engineering Subject \texttt{SUBJ-000001}.

\begin{evidencebox}
canonical_subject_id: SUBJ-000001
canonical_label:      Microscope Workstation
subject_form:         entity
identity_status:      resolved
mention_ids:          [MNT-000001]
\end{evidencebox}

\noindent\textit{Originalartefakt:}
\texttt{canonical\_subject\_set.json} aus
\texttt{RUN-000001/ATT-000002}.

\subsubsection{Persona-Interpretationen und Consensus/Variance}

Die drei Persona-Interpretationen referenzieren dasselbe Canonical Subject.
Für die Dimension \texttt{information\_type} lagen drei unterschiedliche
Klassifikationen vor. Der Consensus-Zustand wurde deshalb mit
\texttt{divergent} und \texttt{review\_attention\_required=true} persistiert.

\begin{evidencebox}
canonical_subject_id: SUBJ-000001

information_type:
  confidence:                low
  consensus_level:           divergent
  review_attention_required: true

  PERSONA_LEGACY_LITERAL_INTERPRETER:
    function

  PERSONA_LEGACY_SYSTEMS_ENGINEERING_INTERPRETER:
    logical_element

  PERSONA_LEGACY_SKEPTICAL_AMBIGUITY_INTERPRETER:
    physical_element

uncertainty:
  The exact system boundary and whether this is hardware,
  software, or a combined workstation are not specified.
\end{evidencebox}

\noindent\textit{Originalartefakt:}
\texttt{runs/RUN-000001/artifacts/consensus\_reports/agentic\_ingestion/}
\texttt{ATT-000002/subject\_consensus.json}.

Der Consensus-Zustand enthält an dieser Stelle noch keine autorisierende
Engineering-Entscheidung. Er bündelt die auf denselben Subject-Bezug
zurückgeführten Persona-Ergebnisse und markiert die Dimension, für die Human
Review erforderlich ist.

\subsubsection{Human Engineering Review}

Das Canonical Subject wurde als \texttt{RIT-000016} in das Human Engineering
Review überführt. Der finalisierte Report enthält den freigegebenen Inhalt,
die Human Classification, den Review Outcome sowie die zugehörige Rationale.

\begin{evidencebox}
RIT-000016 -- Microscope Workstation

Review Outcome:
  accepted_with_modification

Reviewed Content:
  The Microscope Workstation is described as providing local session control.

Information Type:
  logical_element

Modality:
  descriptive

Epistemic Status:
  explicit

Human Rationale:
  Logical
\end{evidencebox}

\noindent\textit{Originalartefakt:}
\texttt{reviews/RVD-000004/versions/RVV-000001/finalized/reviewed\_report.md}.

Die zugehörige Revision \texttt{RVR-000009} persistiert zusätzlich die
Dimensionen der Human Decision, die Reviewer-Identität und den Zeitpunkt der
Auswahl.

\begin{evidencebox}
dimension: classification
selected_values:
  - logical_element
  - descriptive
  - explicit
value_origin: item_override
rationale:    Logical
selected_by:  MZ
selected_at:  2026-09-11T05:30:09Z

review_outcome:
  accepted_with_modification
\end{evidencebox}

\noindent\textit{Originalartefakt:}
\texttt{reviews/RVD-000004/versions/RVV-000001/revisions/RVR-000009.json}.

\subsubsection{Approved Engineering Information und Project Fit}

Nach Finalisierung des Reviews wurde der autorisierte Zustand als
\texttt{AIN-000011} persistiert. Das Approved Input bindet den fachlichen
Inhalt an den exakten Review-, Source-, Run- und Attempt-Zustand.

\begin{evidencebox}
approved_input_id:      AIN-000011
authority_state:        active
review_document_id:     RVD-000004
review_version_id:      RVV-000001
review_revision_id:     RVR-000009
review_item_id:         RIT-000016
source_id:              SRC-000005
processing_run_id:      RUN-000001
attempt_id:             ATT-000002

canonical_content:
  title:                Microscope Workstation
  primary_text:         The Microscope Workstation is described as providing local session control.
  information_type:     logical_element
  modality:             descriptive
  epistemic_status:     explicit
\end{evidencebox}

\noindent\textit{Originalartefakt:}
\texttt{approved\_inputs/manifests/AIN-000011.json}.

Für \texttt{SRC-000005} wurde anschließend ein eigener Project-Fit-Zustand
persistiert. Er referenziert weiterhin Source, Source Projection, Run und
Attempt und hält das Ergebnis sowie die verwendete Begründung fest.

\begin{evidencebox}
source_id:         SRC-000005
processing_run_id: RUN-000001
attempt_id:        ATT-000002
outcome:           plausible_in_scope

rationale:
The candidate source describes a Microscope Workstation with local session
control and a Remote Client that issues remote viewing and control requests,
which aligns with the project context of iO Stream for remote collaboration
in digital microscopy. It matches the local/remote user split and the remote
control workflow described in the context sources.
\end{evidencebox}

\noindent\textit{Originalartefakt:}
\texttt{project\_reconciliation/project\_fit/}
\texttt{0c4839947bd3dcaa17d43734c129356542bf46ac2a1d6a60eb52ad070ca8a49d.json}.

\subsubsection{Model Placement und Internal Engineering Model}

Für \texttt{AIN-000011} wurde mit \texttt{MPD-000016} eine explizite Model
Placement Decision gespeichert. Sie legt den Zielbereich und die anzuwendende
Placement Rule fest.

\begin{evidencebox}
decision_id:       MPD-000016
approved_input_id: AIN-000011
outcome:           accepted
selected_rule_id:  ELEMENT_SYSTEM_LOGICAL
reviewer_identity: MZ
rationale:         Systemlevel
reviewed_at:       2026-09-11T09:08:58Z
\end{evidencebox}

\noindent\textit{Originalartefakt:}
\texttt{model\_placement\_reviews/.../decisions/MPD-000016.json}.

Der daraus aufgebaute interne Modellzustand \texttt{IEM-000001} enthält den
Gegenstand als \texttt{IME-000011}. Die Placement Authority bleibt direkt am
Modellelement referenziert.

\begin{evidencebox}
internal_model_element_id: IME-000011
approved_input_id:         AIN-000011
name:                      Microscope Workstation
description:               The Microscope Workstation is described as providing local session control.
element_type:              logical_component
framework_assignment:      FW_SYSTEM_LOGICAL
model_area:                system.logical

placement_authority:
  authority_id:            MPD-000016
  authority_type:          model_placement_decision
\end{evidencebox}

\noindent\textit{Originalartefakt:}
\texttt{internal\_models\_v2/IEM-000001/snapshot.json}.

\subsubsection{Model Quality Decision und Successor IEM}

Im Model-Quality-Schritt wurde für dasselbe interne Modellelement eine
modellbezogene Formulierung vorgeschlagen und durch \texttt{MZ} freigegeben.
Die Entscheidung wurde als \texttt{MQD-000008} persistiert.

\begin{evidencebox}
decision_id:              MQD-000008
internal_model_element_id: IME-000011
decision:                 approved
reviewer_identity:        MZ

approved_name:
  Microscope Workstation

approved_description:
  Logical component described as providing local session control.

rationale:
  Batch-confirmed after the refinement quality contract reported preserved
  meaning, no unsupported information and no explicit Human-attention flag.
\end{evidencebox}

\noindent\textit{Originalartefakt:}
\texttt{model\_quality/decisions/MQD-000008.json}.

Die freigegebene Formulierung wird nicht in den Vorgängerzustand
\texttt{IEM-000001} zurückgeschrieben. Sie erscheint im nachfolgenden
\texttt{IEM-000002} unter derselben internen Elementidentität.

\begin{evidencebox}
internal_model_element_id: IME-000011
name:                      Microscope Workstation
description:               Logical component described as providing local session control.
element_type:              logical_component
framework_assignment:      FW_SYSTEM_LOGICAL
model_area:                system.logical
placement_authority:       MPD-000016
\end{evidencebox}

\noindent\textit{Originalartefakt:}
\texttt{internal\_models\_v2/IEM-000002/snapshot.json}.

\subsubsection{Generiertes SysML v2 und Output-Traceability}

Die deterministische Generierung überführt \texttt{IME-000011} in das finale
SysML-v2-Artefakt. Der relevante Originalausschnitt lautet:

\begin{evidencebox}
part IME_000011 {
    doc /* Engineering name: Microscope Workstation
    Description:
    Logical component described as providing local session control. */
}
\end{evidencebox}

\noindent\textit{Originalartefakt:}
\texttt{data/output/930444/OUT-000001/generated\_model.sysml},
Zeilen~73--77.

Das zugehörige Output-Traceability-Artefakt bindet das erzeugte Symbol an das
Approved Input, die Model Placement Authority, den Successor-IEM-Zustand und
die genaue Position im generierten Artefakt.

\begin{evidencebox}
approved_input_id:                    AIN-000011
authority_id:                         MPD-000016
authority_type:                       model_placement_decision
generated_symbol_id:                  IME_000011
generated_unit_id:                    GSU-000001
source_internal_engineering_model_id: IEM-000002
source_internal_model_element_id:      IME-000011
generated_location:
  start_line: 73
  end_line:   77
\end{evidencebox}

\noindent\textit{Originalartefakt:}
\texttt{data/output/930444/OUT-000001/traceability.json}.

\subsubsection{Technische Validierung und Human Release}

Der publizierte Output enthält einen menschenlesbaren Validation Report. Da der
Report kompakt ist, wird er hier vollständig wiedergegeben.

\begin{evidencebox}
# SysML v2 Validation Report

- Project: `930444`
- Source IEM: `IEM-000002`
- Validation status: **valid**
- Publication gate: **passed**
- Artifact set fingerprint: `217e89b7e5c37a39842f0479076e3389a0c4b8de0847706b6d84cd77a9ff730c`
- Validation result fingerprint: `aa196bed0529b592c28b3b61d18924b31167d56571f9837f900f6ae48deb4fae`
- Validation profile: `TURING_SYSML_V2_VALIDATION 1.0.0`
- Findings: 0

## External validator evidence

- `SYSIDE_CLI` -- SYSIDE Modeler CLI syside 0.10.3
  (b6e216cb48b5336ea48283e99c68a0e10e17b8cc);
  execution=completed; diagnostics=0

## Findings

No validation findings.
\end{evidencebox}

\noindent\textit{Originalartefakt:}
\texttt{data/output/930444/OUT-000001/validation\_report.md}.

Nach der technischen Validierung wurde die Human Release Decision separat
persistiert. Die Entscheidung verweist auf die exakte Final-Model-Review-
Revision, den generierten Artifact-Set-Fingerprint sowie den zugehörigen
Validation Result Fingerprint.

\begin{evidencebox}
final_model_review_decision_id: FRD-000001
decision:                       approved_for_publication
reviewer_identity:              MZ
reviewed_at:                    2026-09-11T09:24:10Z
rationale:                      null

target:
  final_model_review_id:          FMR-000001
  final_model_review_revision_id: FRV-000001
  validation_status:              valid
  publication_gate:               passed
  generated_artifact_set_fingerprint:
    217e89b7e5c37a39842f0479076e3389a0c4b8de0847706b6d84cd77a9ff730c
  validation_result_fingerprint:
    aa196bed0529b592c28b3b61d18924b31167d56571f9837f900f6ae48deb4fae
\end{evidencebox}

\noindent\textit{Originalartefakt:}
\texttt{final\_model\_reviews/FMR-000001/decisions/FRD-000001.json}.

\subsection{Ergänzende Beispiele für weitere persistierte Entscheidungsarten}
\label{app:traceability_930444_additional}

Der Hauptpfad \emph{Microscope Workstation} deckt nicht jede im Referenzlauf
auftretende Entscheidungsart ab. Zwei kurze ergänzende Beispiele dokumentieren
daher eine Relationship Decision sowie eine explizite Target-Model-Formulation
Decision. Diese Beispiele werden nicht als Fortsetzung desselben
Informationspfads behandelt.

\subsubsection{Relationship Decision: Ablehnung und separate akzeptierte Gegenrichtung}

Im Review der System Requirements wurden zwei getrennte Relationship Proposals
für die Subjects \emph{system} und \emph{indicate the current controller}
entschieden. \texttt{SRD-000001} verwirft die Richtung vom Verhalten zum
System. Die eingegebene Live-Demo-Rationale lautete \texttt{vice versa}.

\begin{evidencebox}
decision_id:       SRD-000001
source_subject_id: SUBJ-000002   # indicate the current controller
relationship_kind: performs
target_subject_id: SUBJ-000001   # system
outcome:           rejected
rationale:         vice versa
reviewer_identity: MZ
\end{evidencebox}

\noindent\textit{Originalartefakt:}
\path{subject_review_relationship_decisions/RVD-000003/RVV-000001/SRD-000001.json}.

Eine separate Proposition mit entgegengesetzter Source-/Target-Richtung wurde
als \texttt{SRD-000012} akzeptiert. Es handelt sich nicht um eine Revision von
\texttt{SRD-000001}, sondern um eine eigenständige Relationship Decision.

\begin{evidencebox}
decision_id:       SRD-000012
source_subject_id: SUBJ-000001   # system
relationship_kind: performs
target_subject_id: SUBJ-000002   # indicate the current controller
outcome:           accepted
rationale:         null
reviewer_identity: MZ
\end{evidencebox}

\noindent\textit{Originalartefakt:}
\path{subject_review_relationship_decisions/RVD-000003/RVV-000001/SRD-000012.json}.

\subsubsection{Target-Model-Formulation: Microscope Operator}

Für das interne Modellelement \texttt{IME-000001} wurde eine separate
Target-Model-Formulation Decision persistiert. Der ausgewählte Zustand legt
fest, dass der Gegenstand formal materialisiert wird, und bindet die konkrete
Formulierung an einen Target-Model-Pattern- und Target-Notation-Construct-Bezug.

\begin{evidencebox}
decision_id:                       TFD-000001
authority_subject_id:              IME-000001
subject_kind:                      element
selected_relevance_outcome:        materialize_formally
selected_target_model_pattern_id:  standalone_reusable_stakeholder_role_part_definition
selected_target_notation_construct_id: TN_003
selected_formulation_text:         part def 'Microscope Operator';
reviewer_identity:                 MZ

rationale:
  Batch-confirmed during deterministic Target-Model preparation from the exact
  pinned Target Model, Target Notation and generation-profile evidence.
\end{evidencebox}

\noindent\textit{Originalartefakt:}
\texttt{target\_model\_formulation/decisions/TFD-000001.json}.

Im generierten SysML-v2-Artefakt erscheint derselbe Engineering-Gegenstand als
Symbol \texttt{IME\_000001}:

\begin{evidencebox}
part def IME_000001 {
    doc /* Engineering name: Microscope Operator
    Description:
    Local person working at the microscope workstation and responsible for the local session. */
}
\end{evidencebox}

\subsection{Umfang des im Anhang wiedergegebenen Evidenzsatzes}
\label{app:traceability_930444_scope}

Die vorstehenden Auszüge bilden bewusst keinen vollständigen Dump des
Runtime-Verzeichnisses. Vollständig wiedergegeben werden die kurzen, für die Evaluation
aufbereiteten Engineering-Eingangsdokumente und der kompakte Validation Report. Bei den
umfangreicheren Maschinenartefakten werden diejenigen Felder gezeigt, die zur
Rekonstruktion des exemplarischen Pfads und der ergänzenden Entscheidungsarten
erforderlich sind. Die Originalartefakte bleiben jeweils über ihre eindeutigen
IDs, Dateipfade und persistierten Fingerprints identifizierbar.

Der Anhang ergänzt damit die Ergebnisdarstellung in
Abschnitt~\ref{sec:results_traceability}. Die Einordnung, welche Bedeutung die
persistierten Zwischenzustände, Human Decisions und Authority-Bindungen für die
Kontrollierbarkeit des KI-gestützten Engineering-Prozesses besitzen, erfolgt
getrennt in Kapitel~\ref{sec:discussion}.
